\chapter{Arquitetura de Software}

\section{Introdução}

Todo sistema é desenvolvido com um propósito. A forma de saber se o sistema está atendendo ao seu propósito é estabelecendo requisitos que impõem restrições e direcionam o processo de desenvolvimento para a direção correta. Os requisitos em si são uma grande área de pesquisa, mas podem ser principalmente divididos em dois grupos: funcionais e não funcionais.

\section{Requisitos Funcionais e Não Funcionais}

Requisitos funcionais são restrições que especificam o que o sistema deve ser capaz de fazer. Um requisito funcional para uma rede social pode ser a capacidade de postar fotos. Por outro lado, requisitos não funcionais são restrições que definem como algo deve ser feito. Um requisito não funcional para a rede social mencionada pode ser o armazenamento seguro das fotos dos usuários, impedindo acesso não autorizado a elas.

Em resumo, requisitos funcionais definem funcionalidades, enquanto os não funcionais estabelecem expectativas de qualidade.

\section{Definição de Arquitetura de Software}

A arquitetura de software é definida como uma coleção de decisões de design que afetam a estrutura, comportamento e qualidade geral de um sistema de software, servindo como base para decisões subsequentes. É o campo de pesquisa que lança luz sobre requisitos não funcionais. Ao longo dos anos, muitas dessas decisões de design de qualidade foram agrupadas no que são chamados de padrões arquiteturais ou simplesmente padrões.

\section{Padrões Arquiteturais}

Padrões podem ser definidos como conceitos reutilizáveis que podem ser aplicados para resolver questões similares em sistemas em contextos similares, considerando necessidades não funcionais. Isso significa que um padrão pode ser aplicado a qualquer sistema, independentemente de suas funcionalidades, desde que tenham requisitos de qualidade similares.

Se funcionalidades fossem a única parte importante, não haveria necessidade de ter arquitetura de software. Mas ao mesmo tempo, isso não é o que acontece na realidade. Requisitos de qualidade acabam não sendo tão relevantes quanto deveriam ao escolher padrões arquiteturais para um sistema na indústria, e isso traz consequências de longo prazo para confiabilidade, manutenibilidade, testabilidade, entre outras complicações.

\section{Princípios SOLID}

Os princípios SOLID são um conjunto de diretrizes para o design de software orientado a objetos que visam criar sistemas mais flexíveis, manuteníveis e extensíveis. Estes princípios são:

\begin{itemize}
    \item \textbf{Single Responsibility Principle (SRP)}: Uma classe deve ter apenas uma razão para mudar, ou seja, deve ter apenas uma responsabilidade.
    
    \item \textbf{Open/Closed Principle (OCP)}: Entidades de software devem estar abertas para extensão, mas fechadas para modificação.
    
    \item \textbf{Liskov Substitution Principle (LSP)}: Objetos de uma superclasse devem ser substituíveis por objetos de suas subclasses sem alterar a funcionalidade do programa.
    
    \item \textbf{Interface Segregation Principle (ISP)}: Clientes não devem ser forçados a depender de interfaces que não utilizam.
    
    \item \textbf{Dependency Inversion Principle (DIP)}: Módulos de alto nível não devem depender de módulos de baixo nível. Ambos devem depender de abstrações.
\end{itemize}

Estes princípios são especialmente relevantes para o Modelo de Atores, pois cada ator pode ser visto como uma unidade que segue o princípio de responsabilidade única, enquanto a comunicação baseada em mensagens facilita a inversão de dependências e a segregação de interfaces.

\section{Avaliação de Adequação de Padrões}

A escolha de padrões arquiteturais é crucial para o sucesso de um projeto de software. Na prática, a funcionalidade é considerada muito mais frequentemente do que a qualidade ao selecionar padrões arquiteturais. Este trabalho avalia se o Modelo de Atores pode elevar a qualidade em outros tipos de domínios.

Para avaliar a adequação de um padrão arquitetural, é necessário considerar:

\begin{itemize}
    \item \textbf{Requisitos de Qualidade}: Como o padrão atende aos requisitos não funcionais do sistema?
    
    \item \textbf{Contexto de Aplicação}: O padrão é adequado para o domínio e ambiente de execução?
    
    \item \textbf{Trade-offs}: Quais são os benefícios e custos de aplicar o padrão?
    
    \item \textbf{Manutenibilidade}: Como o padrão afeta a facilidade de manutenção e evolução do sistema?
\end{itemize}

\section{Justificativa para o Modelo de Atores}

O propósito deste trabalho é experimentar com um padrão arquitetural específico chamado Modelo de Atores. A ideia é usar este padrão em um projeto que está bem distante dos requisitos de qualidade que ele tradicionalmente atende e entender quais são as consequências.

Esta abordagem experimental permitirá:

\begin{itemize}
    \item Explorar os limites do Modelo de Atores em contextos não distribuídos
    \item Avaliar se os benefícios do padrão se mantêm em aplicações centralizadas
    \item Identificar possíveis adaptações necessárias para diferentes domínios
    \item Contribuir para o conhecimento sobre a aplicabilidade de padrões arquiteturais
\end{itemize}

